Geografická mapa podílu top 10 \% na celkovém bohatství v~\ac{EU} ukazuje výrazné kontrasty: zatímco skandinávské státy a~Belgie udržují podíl top 10 \% relativně nižší (díky vysoké odborové hustotě, erga omnes a~progresivnímu zdanění majetku), Litva, Lotyšsko nebo i~nové členské státy \ac{EU} vykazují extrémní majetkovou koncentraci.
ČR se nachází v~středním poli, avšak s~rostoucím trendem (srov.~obr.~\ref{fig:wealth_top10_timeline}): stojí tedy na rozcestí mezi trajektorií baltských zemí (vysoká koncentrace bez kompenzačních institucí) a~trajektorií kontinentální Evropy (umírněná koncentrace se silnou \ac{KS}).
Mezinárodní výzkum konzistentně ukazuje, že silnější kolektivní vyjednávání koreluje s~nižší majetkovou koncentrací prostřednictvím dvou mechanismů: vyšší mzdový podíl na \ac{HDP} zvyšuje schopnost spodních decilů akumulovat majetek, a~silnější \ac{KS} snižují pravděpodobnost švarcsystému a~nezdaněných pracovních odměn pro top manažery.
Politická doporučení: kombinace erga omnes s~transparentnějšími výplatními tabulkami v~\ac{KS} by snížila nepřezřejmý prostor pro mimomzdové benefity koncentrované v~horních příjmových skupinách.

\acs{var-g-M}
Geografická mapa podílu top 10 \% na celkovém bohatství v~\ac{EU} ukazuje výrazné kontrasty: zatímco skandinávské státy a~Belgie udržují podíl top 10 \% relativně nižší (díky vysoké odborové hustotě, erga omnes a~progresivnímu zdanění majetku), Litva, Lotyšsko nebo i~nové členské státy \ac{EU} vykazují extrémní majetkovou koncentraci.
ČR se nachází v~středním poli, avšak s~rostoucím trendem (srov.~obr.~\ref{fig:wealth_top10_timeline}): stojí tedy na rozcestí mezi trajektorií baltských zemí (vysoká koncentrace bez kompenzačních institucí) a~trajektorií kontinentální Evropy (umírněná koncentrace se silnou \ac{KS}).
Mezinárodní výzkum konzistentně ukazuje, že silnější kolektivní vyjednávání koreluje s~nižší majetkovou koncentrací prostřednictvím dvou mechanismů: vyšší mzdový podíl na \ac{HDP} zvyšuje schopnost spodních decilů akumulovat majetek, a~silnější \ac{KS} snižují pravděpodobnost švarcsystému a~nezdaněných pracovních odměn pro top manažery.
Politická doporučení: kombinace erga omnes s~transparentnějšími výplatními tabulkami v~\ac{KS} by snížila nepřezřejmý prostor pro mimomzdové benefity koncentrované v~horních příjmových skupinách.

\acs{var-g-M}
Geografická mapa podílu top 10 \% na celkovém bohatství v~\ac{EU} ukazuje výrazné kontrasty: zatímco skandinávské státy a~Belgie udržují podíl top 10 \% relativně nižší (díky vysoké odborové hustotě, erga omnes a~progresivnímu zdanění majetku), Litva, Lotyšsko nebo i~nové členské státy \ac{EU} vykazují extrémní majetkovou koncentraci.
ČR se nachází v~středním poli, avšak s~rostoucím trendem (srov.~obr.~\ref{fig:wealth_top10_timeline}): stojí tedy na rozcestí mezi trajektorií baltských zemí (vysoká koncentrace bez kompenzačních institucí) a~trajektorií kontinentální Evropy (umírněná koncentrace se silnou \ac{KS}).
Mezinárodní výzkum konzistentně ukazuje, že silnější kolektivní vyjednávání koreluje s~nižší majetkovou koncentrací prostřednictvím dvou mechanismů: vyšší mzdový podíl na \ac{HDP} zvyšuje schopnost spodních decilů akumulovat majetek, a~silnější \ac{KS} snižují pravděpodobnost švarcsystému a~nezdaněných pracovních odměn pro top manažery.
Politická doporučení: kombinace erga omnes s~transparentnějšími výplatními tabulkami v~\ac{KS} by snížila nepřezřejmý prostor pro mimomzdové benefity koncentrované v~horních příjmových skupinách.

\acs{var-g-M}
Geografická mapa podílu top 10 \% na celkovém bohatství v~\ac{EU} ukazuje výrazné kontrasty: zatímco skandinávské státy a~Belgie udržují podíl top 10 \% relativně nižší (díky vysoké odborové hustotě, erga omnes a~progresivnímu zdanění majetku), Litva, Lotyšsko nebo i~nové členské státy \ac{EU} vykazují extrémní majetkovou koncentraci.
ČR se nachází v~středním poli, avšak s~rostoucím trendem (srov.~obr.~\ref{fig:wealth_top10_timeline}): stojí tedy na rozcestí mezi trajektorií baltských zemí (vysoká koncentrace bez kompenzačních institucí) a~trajektorií kontinentální Evropy (umírněná koncentrace se silnou \ac{KS}).
Mezinárodní výzkum konzistentně ukazuje, že silnější kolektivní vyjednávání koreluje s~nižší majetkovou koncentrací prostřednictvím dvou mechanismů: vyšší mzdový podíl na \ac{HDP} zvyšuje schopnost spodních decilů akumulovat majetek, a~silnější \ac{KS} snižují pravděpodobnost švarcsystému a~nezdaněných pracovních odměn pro top manažery.
Politická doporučení: kombinace erga omnes s~transparentnějšími výplatními tabulkami v~\ac{KS} by snížila nepřezřejmý prostor pro mimomzdové benefity koncentrované v~horních příjmových skupinách.

\acs{var-g-M}
\input{texparts/python/wealth_top10_map}



